Dominique répond à un QCM comportant 10 questions.
Pour chaque question, il est proposé 4 réponses dont une seule est exacte.
Dominique répond au hasard à chacune des 10 questions en cochant, pour chaque question, exactement une case parmi les 4.
Pour chacune des questions, la probabilité qu'il réponde correctement est donc $\frac{1}{4}$.

On note $X$ la variable aléatoire qui compte le nombre de bonnes réponses à ce QCM.

\begin{enumerate}
\item 
Pour chaque question, il y a deux issues possibles: Dominique donne la bonne réponse, avec une probabilité $p=\frac{1}{4}$, ou Dominique donne une mauvaise réponse.

On effectue cette expérience élémentaire 10 fois et de façon indépendante.

La variable aléatoire $X$ qui compte le nombre de bonnes réponses à ce QCM suit donc la loi binomiale de paramètres $n=10$ et $p=\frac{1}{4}$.

\item 
La probabilité que Dominique obtienne exactement 5 bonnes réponses est :

$P\left (X=5\right ) = \binom{10}{5} \times \left (\frac{1}{4}\right )^5 \times \left (1-\frac{1}{4}\right )^{10-5}\approx \np{0,0584}$

\item 
L'espérance de $X$ est $E(X)=np = 10\times \frac{1}{4}= 2,5$.

En procédant comme Dominique, un élève aura en moyenne $2,5$ bonnes réponses.

\item On suppose dans cette question qu'une bonne réponse rapporte un point et qu'une mauvaise réponse fait perdre $0,5$ point. La note finale peut donc être négative.

On note $Y$ la variable aléatoire qui donne le nombre de points obtenus.

	\begin{enumerate}
		\item 
Pour réaliser l'évènement $(Y=10)$, il faut que toutes les réponses soient justes donc:
$P\left (Y = 10\right )=P\left (X = 10\right ) = \left (\frac{1}{4}\right )^{10}$.		
				
		\item 
\begin{list}{\textbullet}{On peut calculer la valeur de $Y$ en fonction de celle de $X$}
\item On a vu que $(Y = 10)$ équivalait à $(X=10)$.
\item Si $X = 9$, il y a 9 bonnes réponses qui rapportent 9 points, et une mauvaise réponse qui en retire un demi; donc $Y=8,5$.
\item Si $X = 8$, il y a 8 bonnes réponses qui rapportent 8 points, et deux mauvaises réponses qui retirent deux demi-points; donc $Y=7$.
\item Etc.
\end{list}

On regroupe tous les résultats possibles dans un tableau.

\begin{center}
\begin{tabularx}{\linewidth}{| *{12}{>{\centering\arraybackslash} X|}}
\hline
$X$ & 0 & 1 & 2 & 3 & 4 & 5 & 6 & 7 & 8 & 9 & 10\\
\hline
$Y$ & $-5$ & $-3,5$ & $-2$ & $-0,5$ & $1$ & $2,5$ & $4$ & $5,5$ & $7$ & $8,5$ & $10$ \\
\hline
\end{tabularx}
\end{center}

Donc la note finale de Dominique est positive à partir de 4 bonnes réponses.

		\item 
D'après le tableau précédent, $P\left (Y\leqslant 0\right )=P\left (X\leqslant 3\right )$.

Donc $P\left (Y\leqslant 0\right ) \approx 0,78$.

		\item
$X$ prend toutes les valeurs entières entre 0 et 10. 

Pour $X$ bonnes réponses qui rapportent $X$ points, il y en a $10-X$ mauvaises qui retirent $0,5\left (10-X\right )$ points.

Donc $Y=X-0,5\left (10-X\right )=X-5+0,5X=1,5X-5$.

		\item 
D'après la linéarité de l'espérance mathématique, on a:

$E\left (Y\right ) = E\left (1,5X-5\right )=1,5\times E\left (X\right ) - 5 = 1,5\times 2,5 - 5 = -1,25$.		
	\end{enumerate}
\end{enumerate}

\bigskip
